\documentclass[12pt]{article}
\usepackage{graphicx} % Required for inserting images
\usepackage{enumitem}
\setlist[enumerate]{label=(\Alph*)}
\title{Computer Network}
\author{}
\date{}

\begin{document}

\maketitle

\section*{Class Question Notes}
\subsection*{Question 1}
\textbf{Find the class B address from the following}\\
\begin{enumerate}[label=(\Alph*)]
    \item $01111111\cdot01010101\cdot 11111110 \cdot 00001111$
    \item $11101111\cdot01001110\cdot11001100\cdot01010011$
    \item $10001111\cdot00000011\cdot11111100\cdot00111100$
    \item $11011111\cdot11001111\cdot11100010\cdot11111010$
\end{enumerate}


\subsection*{Question 2}
\textbf{Find the invalid IP address from the following choices?(Assuming Class-full addressing scheme is followed)} \\
\begin{enumerate}[label=(\Alph*)]
    \item $150.168.10.1$
    \item $190.100.1.100$
    \item $10.256100.100$
    \item $80.10.254.100$
\end{enumerate}

\subsection*{Question 3}
\textbf{Which of the following addresses can be used for interprocess
communication in a host?}
\begin{enumerate}[label=(\Alph*)]
    \item $192.168.100.100$
    \item $127.100.100.100$
    \item $10.100.100.100$
    \item $172.16.100.100$
\end{enumerate}

\subsection*{Question 4}
\textbf{The Dotted decimal notation (DDN) format for the given
Hexadecimal notation (HDN) $C22F1582$ is}
\begin{enumerate}[label=(\Alph*)]
    \item $194.50.21.145$
    \item $194.47.21.130$
    \item $194.45.21.120$
    \item $194.47.20.130$
\end{enumerate}

\subsection*{Question 5}
\textbf{The Dotted decimal notation (DDN) format for the given
Hexadecimal notation (HDN) $172A84C8$}
\begin{enumerate}[label=(\Alph*)]
    \item $24.40.132.200$
    \item $23.42.132.200$
    \item $23.42.130.200$
    \item $23.42.132.198$
\end{enumerate}

\subsection*{Question 6}
\textbf{Suppose, instead of using 16 bits for network part of a Class B, 20
bits had been used. Then the number of Class B networks and hosts
per network are}
\begin{enumerate}[label=(\Alph*)]
    \item $2^{10}, 2^{12}$
    \item $2^{18}, 2^{12}$
    \item $2^{18}, 2^{12} - 2$
    \item $2^{10}, 2^{12} - 2$
\end{enumerate}

\subsection*{Question 7}
\textbf{Number of Networks and Number of Host in class B are $2m$,$(2n - 2)$ respectively. Then the relation between $m$ and $n$ is}
\begin{enumerate}[label=(\Alph*)]
    \item $3m=2m$
    \item $7m=8n$
    \item $8m=7n$
    \item $2m=3n$
\end{enumerate}

\subsection*{Question 8}
\textbf{How many networks are possible in a class B addressing system ?
(Assuming Classful addressing scheme is followed.)}
\begin{enumerate}[label=(\Alph*)]
    \item $2^{16}$
    \item $2^{14}$
    \item $2^8-2$
    \item $2^{16} - 2$
\end{enumerate}

\subsection*{Question 9}
\textbf{How many hosts can be present in a class C network ? (Assuming
Classful addressing scheme is followed.)}
\begin{enumerate}[label=(\Alph*)]
    \item $2^{21}$
    \item ${2^{21}- 1}$
    \item $2^{16}$
    \item $2^8-2$
\end{enumerate}

\subsection*{Question 10}
\textbf{How many bits are allocated for NID and HID in $23.192.157.234$
address ? (Assuming Classful addressing scheme is followed.)}
\begin{enumerate}[label=(\Alph*)]
    \item $16, 16$
    \item $8,16$
    \item $8,24$
    \item $24, 8$
\end{enumerate}

\subsection*{Question 11}
\textbf{In classful addressing, a large part of the available addresses are}\
\begin{enumerate}[label=(\Alph*)]
    \item Dispersed
    \item Blocked
    \item Blocked
    \item Reserved
\end{enumerate}

\subsection*{Question 12}
\textbf{What is the possible number of networks and addresses in each
network under class B addresses in IPv4 addressing format.}
\begin{enumerate}[label=(\Alph*)]
    \item $2^{16}, 2^{16}$
    \item $2^{16}, 2^{16} - 2$
    \item $2^{14}, 2^{16} - 2$
    \item $2^{14}, 2^{16}$
\end{enumerate}

\subsection*{Question 13}
\textbf{IP Address $200.198.32.65$ belong to which class ?}
\begin{enumerate}[label=(\Alph*)]
    \item Class A
    \item Class B
    \item Class C
    \item Class D
\end{enumerate}

\subsection*{Question 14}
\textbf{Percent of Addresses occupied by Class D ?}
\begin{enumerate}[label=(\Alph*)]
    \item $50\%$
    \item $25\%$
    \item $6.25\%$
    \item $12.5\%$
\end{enumerate}

\subsection*{Question 15}
\textbf{In IPv4 addressing format, the number of networks all allowed
under class C addresses is}
\begin{enumerate}[label=(\Alph*)]
    \item $2^{24}$
    \item $2^{7}$
    \item $2^{14}$
    \item $2^{21}$
\end{enumerate}

\subsection*{Question 16}
\textbf{A host with IP address $10.100.100.100$ wants to use loopback testing. What are the source and destination addresses ? (Assuming Classful addressing scheme is followed.)}
\begin{enumerate}[label=(\Alph*)]
    \item $10.100.100.100$ and $10.100.100.100$
    \item $10.100.100.100$ and $255.255.255.255$
    \item $10.100.100.100$ and $127.1.100.1$
    \item $127.100.100.100$ and $10.100.100.100$
\end{enumerate}

\subsection*{Question 17}
\textbf{$100.86.95.75$, $157.192.190.253$, $,200.1.56.97$, $10.34.87.95$ Which of the following is common for all
these IP Addresses.}
\begin{enumerate}
    \item Class of IP address
    \item Limited broadcast address
    \item Network address
    \item Direct broadcast address
\end{enumerate}

\subsection*{Question 18}
\textbf{For the IP Addresses $132.54.78.98$ identify the Class ,and Limited broadcast Address}
\begin{enumerate}
    \item IP address belong to class A, Limited broadcast address $= 255.255.255.255$
    \item IP address belong to class B, Limited broadcast address $= 130.255.255.255$
    \item IP address belong to class B, Limited broadcast address $= 255.255.255.255$
    \item IP address belong to class A, Limited broadcast address $= 130.54.255.255$
\end{enumerate}

\subsection*{Question 19}
\textbf{One host having IP address 200.187.96.0, sends a message to a host with IP address 205.54.83.97, what will be the destination address attached to message by source?}
\begin{enumerate}
    \item $205.54.83.97$
    \item $205.54.83.255$
    \item $205.54.83.0$
    \item Not possible
\end{enumerate}

\subsection*{Question 20}
\textbf{Which of the following can be used as a source IP as well as destination IP ?}
\begin{enumerate}
    \item $23.0.0.97$
    \item $255.255.255.255$
    \item $157.54.255.255$
    \item $15.255.255.255$
\end{enumerate}

\subsection*{Question 21}
\textbf{Which of the following IP address can be given to a computer as a host?}
\begin{enumerate}
    \item $32.0.0.0$
    \item $255.255.255.255$
    \item $157.54.255.254$
    \item $172.15.0.0$
\end{enumerate}

\subsection*{Question 22}
\textbf{If NID = 200.200.200.0 and the subnet Mask = $255.255.255.224$ then identify:}
\begin{enumerate}[label = (\Roman*.)]
    \item Number of bit borrowed from Host-id.
    \item Number of Subnet possible and their subnet id's.
    \item Number of Host/Subnet.
\end{enumerate}

\subsection*{Question 23}
\textbf{If NID = $200.200.200.0$ and the subnet Mask $255.255.255.44$ then identify}
\begin{enumerate}[label = (\Roman*.)]
    \item Number of bit borrowed from Host-id.
    \item Number of subnet possible and their subnet id's
    \item Number of Host/subnet
\end{enumerate}

\subsection*{Question 24}
\textbf{If NID = $200.200.200.0$ and the subnet Mask = $255.255.255.200$ then identify}
\begin{enumerate}[label = (\Roman*.)]
    \item Number of bit borrowed from Host-id
    \item Number of subnet possible and their subnet id's
    \item Number of Host/subnet
\end{enumerate}

\subsection*{Question 25}
\textbf{If NID = $173.173.0.0$ and the subnet Mask $255.255.128.128$ then identify}
\begin{enumerate}[label = (\Roman*.)]
    \item Number of bit borrowed from Host-id
    \item Number of subnet possible and their subnet id's
    \item Number of Host/subnet
\end{enumerate}

\subsection*{Question 26}
\textbf{If NID = $173.173.0.0$ and the subnet Mask $255.255.255.0$ then identify}
\begin{enumerate}[label = (\Roman*.)]
    \item Number of bit borrowed from Host-id
    \item Number of subnet possible and their subnet id's
    \item Number of Host/subnet
\end{enumerate}

\subsection*{Question 27}
\textbf{Which of the following is the default mask for the address 198.0.46.201? (Assuming Classful addressing scheme is followed)}
\begin{enumerate}
    \item $255.0.0.0$
    \item $255.255.255.0$
    \item $255.255.0$
    \item $255.255.255.255$
\end{enumerate}

\subsection*{Question 28}
\textbf{If a class B network on the Internet has a subnet mask of 255.255.248.0. What is the maximum number of hosts per subnet? (Assuming Classful addressing scheme is followed)}
\begin{enumerate}
    \item 1022
    \item 1023
    \item 2046
    \item 2047
\end{enumerate}

\subsection*{Question 29}
\textbf{A subnet has assigned a subnet mask of 255.255.255.192. What is the maximum number of hosts that can belong to this subnet ?}
\begin{enumerate}
    \item 14
    \item 30
    \item 62
    \item 126
\end{enumerate}

\subsection*{Question 30}
\textbf{In a class B network on the Internet has a subnet mask of 255.255.240.0. What is the maximum number of hosts per subnet? \(Assuming Classful addressing scheme is followed\)}
\begin{enumerate}
    \item 4096
    \item 4096
    \item 4092
    \item 4090
\end{enumerate}

\subsection*{Question 31}
\textbf{An organization has a class B network and wishes to form subnets for 64 departments. The subnet mask would be:}
\begin{enumerate}
    \item 255.255.0.0
    \item 255.255.64.0
    \item 255.255.128.0
    \item 255.255.252.0
\end{enumerate}

\subsection*{Question 32}
\textbf{Consider default subnet mask for a network is  255.255.255.0. How many number of hosts per subnet possible if 'm' bits are borrowed from Host ID (HID)}
\begin{enumerate}
    \item $2^{HID-m}-2$
    \item $2^{HID}$
    \item $2^{HID}-m$
    \item $2^m$
\end{enumerate}

\subsection*{Question 33}
\textbf{A university has LANs with 100 hosts in each LAN. If it uses class B then the subnet mask in Dotted Decimal Notation is}

\subsection*{Question 34}
\textbf{A university has 150 LANs. Use Class B address and then the subnet mask in Doted Decimal notation is}

\subsection*{Question 35}
\textbf{IP Address = 200.200.200.120 Subnet Mask = 255.255.255.240 then find the SID and HID?}

\subsection*{Question 36}
\textbf{IP Address = 200.200.200.120 subnet Mask = 255.255.255.41 then find the SID and HID?}

\subsection*{Question 37}
\textbf{Find the subnet Address for the Following IP Address: 200.34.22.156 Mask: 255.255.255.240}
\begin{enumerate}
    \item $200.33.22.144$
    \item $200.34.22.143$
    \item $200.34.22.13$
    \item $200.34.22.144$
\end{enumerate}

\subsection*{Question 38}
\textbf{If Subnet Mask is 255.255.255.240 then find}
\begin{enumerate}
    \item Number of IP Address/subnet possible
    \item Number of Host/subnet possible
    \item Number of subnet in class A
    \item Number of subnet in class B
    \item Number of suhnet in class C
\end{enumerate}

\subsection*{Question 39}
\textbf{If Subnet Mask is 255.255.252.0 then find}
\begin{enumerate}
    \item Number of IP Address/subnet possible
    \item Number of Host/subnet possible
    \item Number of subnet in class A
    \item Number of subnet in class B
    \item Number of subnet in class C
\end{enumerate}

\subsection*{Question 40}
\textbf{If Subnet Mask is 255.252.0.0 then find}
\begin{enumerate}
    \item Number of IP Address/subnet possible
    \item Number of Host/subnet possible
    \item Number of subnet in class A
    \item Number of subnet in class B
    \item Number of subnet in class C
\end{enumerate}
\end{document}
